\section{Chap 1: Introduction}
\subsection{Modeling}
\noindent\begin{equation*}
    \boxed{\Dot{x} = f(x)}
\end{equation*}

\textbf{Def}: A system is said to be \textbf{autonomous} or \textbf{time invariant} $\Dot{x} = f(x)$ does not depend explicitly on $t$.

\newpar{}
\textbf{Def}: A point is called an \textbf{equilibrium point} of $\Dot{x} = f(x)$ if $x(\tau) = \Bar{x}$ for some $\tau \implies x(t) = \Bar{x} \quad \forall \; t>\tau$ (or $\Dot{x} = 0$)

Solution to the system of equations $\Dot{x} = Ax + Bu$:
\noindent\begin{equation*}
    \boxed{x(t) = e^{At}\cdot x_0 + \int_0^t e^{A(t-\tau)} Bu(\tau) d\tau}
\end{equation*}\\
List of common non-linearities:
\begin{itemize}
    \item \textbf{Relay}: Appear in modeling mode changes. Comparable to a sign function.
    \item \textbf{Saturation}: Modeling components with hard physical limits. Peak actuator force.
    \item \textbf{Dead Zone}: Friction, sensitivities of sensors and actuators.
    \item \textbf{Quantization}: Discretization when converting analog measurements to digital readings.
\end{itemize}
\vspace{5pt}
\textbf{Def}: A 2D \textbf{vector field} can be drawn by assigning a vector $(f(x_1), f(x_2))$ to each point $(x_1, x_2)$ of the 2D plain.
\newpar{}
\textbf{Def}: The family of all trajectories of a dynamical system is called the \textbf{phase portrait}.

\subsection{Linearization}
Linearization generates a Matrix by partially deriving by the state and plugging in the equilibrium point $\Bar{x}$:
\noindent\begin{equation*}
    \boxed{\Dot{x} = f(x) \approx f(x) \cdot \frac{\partial}{\partial x^T} = \begin{bmatrix}
            \frac{\partial}{\partial x_1} f_1 & \cdots & \frac{\partial}{\partial x_n} f_1 \\
            \vdots                            & \ddots & \vdots                            \\
            \frac{\partial}{\partial x_1} f_m & \cdots & \frac{\partial}{\partial x_n} f_m
        \end{bmatrix}}
\end{equation*}
\subsection{Limit Cycles}
\textbf{Def}: An isolated periodic orbit is called a \textbf{limit cycle}.\\
\textbf{Theorem}: \textbf{Poincare-Bendixson Criterion} Consider the system plane $\Dot{x} = f(x)$ and let $M$ be a closed bounded subset of the plane such that:
\begin{itemize}
    \item $M$ contains no equilibtium points, or contains only one equilibrium point such that the Jacobian matrix at this pint has eigenvalues with positive real parts.
    \item Every trajectory starting in $M$ remains in $M$ for all future time.
\end{itemize}
Then $M$ contains a periodic orbit.\\
\textbf{Theorem}: \textbf{Negative Poincare-Bendixson Criterion} If, on a simply connected region $D$ of the plane, the expression $\frac{\partial f_1}{\partial x_1} + \frac{\partial f_2}{\partial x_2}$ is not identically zero and does not change sign, then the system $\Dot{x} = f(x)$ has no periodic orbits lying entierly in $D$.
